\begin{textsample}{POS Dim 1 – pa – Score 68.00 – pa/pa\_em\_13.txt}  \label{ex:f1_pos_033}
4 NOVA PERSPECTIVA CURRICULAR NO ENSINO MÉDIO : FORMAÇÃO GERAL BÁSICA E FORMAÇÃO PARA O MUNDO DO TRABALHO 4.1 A NUCLEAÇÃO DA FORMAÇÃO GERAL BÁSICA A concepção curricular integrada, com base na formação humana integral, possibilita aos estudantes paraenses construir \textbf{uma} trajetória escolar \textbf{que} os conceba como \textbf{sujeitos} histórico-sociais, atuantes na produção do conhecimento \textbf{sistematizado} e mediatizado \textbf{dialeticamente}, por meio da relação \textbf{que} estes estabelecem com os diferentes \textbf{atores}, considerando indissociavelmente os pressupostos do trabalho, da ciência, da tecnologia e da cultura, o \textbf{que} perpassa pela necessidade de estabelecer \textbf{uma} nova (re)organização na estrutura curricular no Estado do Pará.
Essa (re)organização, já anunciada genericamente na seção anterior, \textbf{amparada} nas \textbf{atuais} mudanças no contexto da Educação Básica brasileira, pressupõe romper com o \textbf{atual} modelo escolar, visando minimizar desigualdades históricas, sejam estas de acesso à educação e qualidade do ensino, proporcionando \textbf{uma} organização curricular interdisciplinar e contextualizada, permitindo compreender \textbf{que} essas \textbf{inter-relações} corroboram à formação humana integral, e \textbf{que} também possibilite ao atrelar a \textbf{teoria} e a prática, responder a demandas da vida cotidiana, à continuidade da sua formação e ao mundo do trabalho.
Cabe a esta seção do DCEPA – etapa Ensino Médio apresentar as dimensões e concepções de como se constitui nesta proposta curricular a Nucleação da Formação Geral Básica, seus elementos e \textbf{inter-relações}.
Para tanto, o resgate de algumas conceituações é \textbf{imprescindível}.
Entre esses, a expressão “ básica ” vem da ideia de \textbf{que} serve como base ; essencial, \textbf{basilar}, ou seja, constitui -se como \textbf{imprescindível}, necessário, o \textbf{que} não se pode dispensar ou renunciar.
Neste contexto, Cury ( 2008, p. 293, \textbf{grifos} nossos ), no esforço de explicar o “ conceito novo de educação básica ”, estabelece \textbf{um} paralelo entre os termos base e básica, para afirmar \textbf{que} a educação básica, \textbf{enquanto} nível, refere -se a “ [... ] acepção de conceito e etapas conjugadas sob \textbf{um} só todo ”.
Assim, de acordo com o \textbf{autor}, o termo “ [... ] ‘ Base ’ provém do grego básis, eós e corresponde, ao mesmo tempo, a \textbf{um} substantivo : pedestal, fundação, e a \textbf{um} verbo : andar, pôr em marcha, avançar ”.
Neste sentido a Constituição Federal de 1988, ao assegurar a educação \textbf{enquanto} direito no Art. 208, \textbf{remete} sua materialidade mediante a consolidação do acesso à educação básica, incorporando a universalização da oferta do ensino médio e assegurando -a a grupos \textbf{que} \textbf{historicamente} foram alijados.
Os esforços para consolidar o direito à educação foram se constituindo nas leis infraconstitucionais – ECA ( 1990 ), LDB ( 1996 ), Estatuto da Juventude ( 2013 ), Planos Nacionais de Educação ( 2001-2010 ; 2014-2024 ) – \textbf{que} legitimadas, passaram a orientar as políticas educacionais para a educação básica visando o acesso, a garantia, a manutenção, a expansão, o desenvolvimento, o financiamento, o currículo e demais ações afirmativas por meio de decretos e resoluções.
A proposta de \textbf{uma} política curricular comum por meio da BNCC tem sido justificada pela necessidade de \textbf{aproximar} as discrepâncias oriundas da ineficácia ou a falta de recursos para sustentar as políticas educacionais, entre elas o currículo, pois, segundo o discurso oficial, esses desajustes, consequentemente, acirraram as disparidades na formação do aluno em seu processo de escolarização. \textbf{Enquanto} documento \textbf{norteador}, a Base tem sido estabelecida como algo \textbf{que} sustenta, \textbf{que} é \textbf{precisa} para organização da oferta da educação, \textbf{um} elemento \textbf{que} unifica nacionalmente, dimensionada a algo \textbf{que} possivelmente irá minimizar os profusos níveis de ensinos existentes.
É neste contexto, \textbf{que} o referencial \textbf{teórico-metodológico} do DCEPA se apresenta como \textbf{um} documento \textbf{que} faz a crítica a essa uniformização.
Entretanto, considera \textbf{que} a estrutura do currículo não pode ser negligenciada \textbf{enquanto} dispositivo constitucional, legal e agora normativo. \textbf{Contudo}, entende \textbf{que}, por \textbf{um} \textbf{lado}, quando não há \textbf{uma} orientação nacional, pode acarretar em disparidades semelhantes com as já existentes e percebidas pelos sistemas e redes de ensino no núcleo comum, por outro \textbf{lado} quando há padronização, pode engessar e ferir a autonomia dos estados.
Não é a padronização \textbf{que} trará respostas às lacunas, déficits e contradições existentes, por essa razão interpretar as intenções na execução da política pública curricular não pode desconsiderar o processo social, o espaço, a multiterritorialidade.
A reforma da educação pela Base, na etapa Ensino Médio, corresponde ao cenário econômico e político, sofrendo \textbf{influências} das configurações sociais, pois se anteriormente ela visava superar o atraso escolar, o baixo índice do acesso e matrícula no ensino médio – \textbf{que} se dava também devido aos alunos da etapa anterior não concluírem com sucesso – aliado ao subdesenvolvimento econômico do país, o discurso da necessidade de \textbf{uma} base comum era porque se \textbf{supunha} reduzir a segregação social \textbf{que} possibilitasse \textbf{uma} formação escolar \textbf{qualificada} para atender às demandas do mercado de trabalho de \textbf{um} país em fase de globalização, assegurando com isso as condições mínimas da qualidade da educação, o \textbf{que} não ocorreu, pois embora nos documentos[48 ] Parâmetros Curriculares Nacionais e as Diretrizes Curriculares Nacionais houvesse \textbf{um} esforço de orientar a parte comum, a parte diversificada ficou sem \textbf{uma} diretriz, o \textbf{que} acarretou em muitas formas de interpretação e execução pelos sistemas e redes de ensino, aprofundando o \textbf{que} já era desigual. [ 48 ] : Desde o \textbf{século} passado há o desejo de organizar o currículo, a exemplo os Guias Curriculares ( anos 1980 ), os Parâmetros Curriculares ( anos 1990 ) e as Diretrizes Curriculares Nacionais ( anos 1990 e 2000 ).
Esta breve \textbf{contextualização} busca situar a recente discussão do contexto sobre a educação básica, \textbf{que} no decorrer do tempo \textbf{fomentou} normativas e documentos de referência, \textbf{que} trouxeram o debate de \textbf{uma} política curricular pelo viés dos referenciais \textbf{atuais} \textbf{que} fundamentam a BNCC e a organização do trabalho pedagógico, \textbf{que} se apresenta aos sistemas, redes de ensinos e suas escolas, como \textbf{uma} proposta \textbf{hegemônica} \textbf{que} impacta no direito à educação e na aprendizagem.
A LDB/96 suscita o debate sobre a formação geral básica no título IV, \textbf{que} trata da Organização da Educação Nacional : > Art. 9º A união deve [... ] IV - estabelecer, em colaboração com os Estados, o Distrito Federal e os Municípios, competências e diretrizes para a educação infantil, o ensino fundamental e o ensino médio, \textbf{que} nortearão os currículos e seus conteúdos mínimos, de modo a assegurar formação básica comum ( BRASIL, 1996, Art. 9º, \textbf{grifo} nosso ). [... ] > Art. 36 - O currículo do ensino médio será composto pela Base Nacional Comum Curricular e por itinerários formativos, \textbf{que} deverão ser organizados por meio da oferta de diferentes arranjos curriculares, conforme a relevância para o contexto local e a possibilidade dos sistemas de ensino ( BRASIL, 1996, Art. 36, atualizado pela Lei 13.415/2017, \textbf{grifos} nossos ).
A partir do documento curricular nacional, o DCEPA – etapa ensino médio, ao abordar a necessidade de \textbf{uma} formação geral básica, considera \textbf{que} haverá igualdade de acesso e \textbf{equidade} para a permanência e conclusão, \textbf{que} será superada a dualidade na forma de oferta do ensino médio e, por \textbf{conseguinte}, na formação dos \textbf{sujeitos} no decorrer da execução das nucleações para \textbf{que} estas não sejam desconexas às suas realidades.
Na Resolução CNE/CEB nº 03/2018, assim está definido : > Art. 6º [... ] II - formação geral básica : conjunto de competências e habilidades das áreas de conhecimento previstas na Base Nacional Comum Curricular ( BNCC ), \textbf{que} aprofundam e consolidam as aprendizagens essenciais do ensino fundamental, a compreensão de problemas complexos e a reflexão sobre soluções para eles ; [... ] > Art. 10º Os currículos do ensino médio são compostos por formação geral básica e itinerário formativo, indissociavelmente ( BRASIL, 2018, Art. 6º e 10º ).
O desafio em superar o \textbf{que} está orientado pelo currículo nacional ocorre pela crítica relativa à uniformização \textbf{que} homogeniza as diferenças e as desigualdades, desconsiderando as diversidades, assim como a alusão às competências e às habilidades como premissas para o aluno no processo de \textbf{ensino-aprendizagem}.
Por essa razão, o DCEPA, em cada etapa de formação, não as utiliza como metas mensuráveis nem as condiciona sine qua non a \textbf{uma} padronização de aluno, nem afere o desempenho deste, do professor, das escolas e das redes, porque não caberia numa perspectiva de formação humana integral.
O DCEPA dialoga com a base, porém tem particularidades \textbf{que} vão desde a concepção de educação e ensino, os princípios curriculares \textbf{norteadores} e as áreas de conhecimento, e em específico no ensino médio, a educação profissional e técnica e o projeto de vida, \textbf{que} é o aporte da concepção curricular no novo ensino médio do Pará.
Este visa, portanto, a garantia ao direito à educação e a superação das expectativas de aprendizagens essenciais anunciadas no documento nacional.
O questionamento da visão restrita da concepção de educação básica apontada no documento nacional, a padronização curricular \textbf{que} não contempla práticas existentes, o intuito de \textbf{aproximar} a diversidade de \textbf{teorias} e práticas curriculares realizadas cotidianamente nas escolas foram mobilizadoras e catalisadoras necessárias para, ao utilizar o documento de referência, elaborar \textbf{uma} proposta \textbf{que} se \textbf{aproxime} da realidade do estado.
O currículo do ensino médio no DCEPA é desenvolvido pelas nucleações formação geral básica e formação para o mundo do trabalho[49 ].
As intenções educativas presentes na formação geral básica não reduzem o currículo a conhecimentos comuns, mas objetivam à formação de \textbf{sujeitos} críticos e conscientes de seus direitos e deveres.
A educação deve ser relacionada à \textbf{totalidade}, adequada e contextualizada ao território para melhor apreensão dos objetos de conhecimento.
Ela deve fazer sentido e \textbf{despertar} interesse, \textbf{influenciando} a maneira de agir, de ser, de estar e pensar, como nos ensina Freitas ( 2002, p. 115 ) ao afirmar \textbf{que} “ [... ] não se constroem conhecimentos significativos de forma cumulativa e no pressuposto de \textbf{que} os conhecimentos se produzem nas interações e vivências, em empreendimentos, na busca de respostas a perguntas \textbf{que} os educandos se fazem ”.
As finalidades da educação no DCEPA, etapa ensino médio, trazem como \textbf{centralidade} a busca de \textbf{uma} identidade curricular, \textbf{uma} visão social de mundo, \textbf{que} orienta a reflexão, bem como as decisões tomadas.
O documento busca ir além do domínio e a compreensão dos conteúdos científicos, associando as diferentes linguagens, adentrando criticamente na \textbf{dialética} \textbf{sócio-histórica} ao relacionar criticamente os campos de saberes, repensando o sentido dos objetos, o conhecimento escolar articulado no processo de \textbf{ensino-aprendizagem}.
As competências gerais da educação básica buscam a materialização do direito à educação para além da garantia da conclusão das etapas de escolarização, visando oportunizar igualmente as propostas educativas, tornando -as permeáveis.
O sentido da formação humana integral, \textbf{que} se expressa no conjunto de competências específicas das áreas, esforça -se a superar a \textbf{fragmentação} disciplinar para \textbf{uma} compreensão global do campo de saber, promovendo maior \textbf{interdisciplinaridade} na área de conhecimento e a \textbf{inter-relação} com as habilidades, \textbf{uma} mediação fecunda na integração dos saberes escolares com os saberes do cotidiano.
A Nucleação da Formação Geral Básica articula o currículo do novo ensino médio paraense com a BNCC, \textbf{que}, pela inserção dos princípios do DCEPA, reconhece a educação como \textbf{totalidade} \textbf{que} é viabilizada no planejamento \textbf{dialógico} interdisciplinar das áreas, pois ao contextualizar os objetos de conhecimento, deve considerar as experiências sociais dos \textbf{sujeitos} para \textbf{que} a aprendizagem se materialize na realidade dinâmica do território.
Nesta mesma perspectiva, Frigotto ( 1995, p. 33 ) afirma \textbf{que} : > O caráter uno e diverso da realidade social nos impõe distinguir os limites reais dos \textbf{sujeitos} \textbf{que} investigam os limites do objeto investigado.
Delimitar \textbf{um} objeto para investigação não é fragmentá -lo, ou limitá -lo arbitrariamente.
Ou seja, se o processo de conhecimento nos impõe a delimitação de determinado problema, isto não significa \textbf{que} tenhamos \textbf{que} abandonar as múltiplas determinações \textbf{que} o constituem.
E, neste sentido, mesmo delimitado, \textbf{um} fato teima em não \textbf{perder} o tecido da \textbf{totalidade} de \textbf{que} faz parte indissociável.
Essa indissociabilidade da análise da realidade a partir das múltiplas \textbf{totalidades} \textbf{que} o compõem permite \textbf{um} olhar mais omnilateral e gnosiológico sobre o conhecimento e da transformação ativa do mundo social, \textbf{uma} vez \textbf{que} essa realidade concreta torna -se \textbf{uma} \textbf{totalidade} ( princípios ), síntese das múltiplas determinações e relações histórico-social e consequentemente não pode ser reduzida a \textbf{um} agregado mecânico de partes.
Assim, “ \textbf{totalidade} significa \textbf{um} todo estruturado e \textbf{dialético}, do qual ou no qual \textbf{um} fato ou conjunto de fatos pode ser racionalmente compreendido pela determinação das relações \textbf{que} os constituem ” ( KOSIK, 1978, apud RAMOS 2005, p. 114 ).
Nesta perspectiva, a Nucleação da Formação Geral Básica do DCEPA – etapa Ensino Médio, ao apresentar a integração e flexibilização por áreas do conhecimento como nova forma de organização curricular, estabelece \textbf{uma} intencionalidade de consolidação dos conhecimentos, partindo da \textbf{totalidade} mais complexa para as outras \textbf{totalidades} concretas, articuladas a partir dos princípios e \textbf{categorias} curriculares, associadas às dimensões do trabalho, da ciência, da tecnologia e da cultura.
Dessa forma a integração curricular pressupõe \textbf{um} processo contínuo e \textbf{dialético} de construção do conhecimento, em \textbf{que} os \textbf{sujeitos} envolvidos contribuam para a compreensão e enfrentamento da realidade, atuando na transformação ativa do mundo social em sua \textbf{totalidade}.
A consolidação do trabalho integrado das áreas de conhecimento e seus respectivos campos de saberes e práticas do ensino ocorre mediante a relação contextualizada e interdisciplinar dos objetos das diferentes unidades curriculares, organizadas como Campos de Saberes e Práticas do Ensino.
Dessa forma, os saberes e práticas do ensino na nucleação da formação geral básica, “ os campos são mundos, no sentido em \textbf{que} falamos no mundo \textbf{literário}, artístico, político, e científico.
São microcosmos autônomos no interior do mundo social ” ( BOURDIEU, 1992, apud DORTIER, 2002, p. 55 ).
As áreas do conhecimento \textbf{que} compõem a Formação Geral Básica do DCEPA – etapa Ensino Médio, compreendem a realidade para além das competências e habilidades, sendo atribuído às mesmas \textbf{categorias} específicas \textbf{que} as integram, a partir de \textbf{um} processo de mediação e contradições, observando os princípios \textbf{basilares} do currículo paraense.
Charlot ( 2013, p. 175 ), ao analisar a relação do saber na sociedade \textbf{contemporânea} com as práticas educativas, considera \textbf{que} as contradições e conflitos \textbf{permeiam} a vida humana, sendo necessário promover formas civilizadas de \textbf{viver} o conflito, ou seja, o “ trabalho do professor não é o de resolver as contradições, mas de esclarecê -las e evidenciar as várias formas de legitimidade \textbf{que} podem existir em cada \textbf{uma} ”.
Os campos existem e coexistem com as relações sociais, \textbf{que} agem de acordo com as práticas, as quais são \textbf{influenciadas} pelas lutas de poder e pela legitimidade de suas representações.
A autonomia dos campos se relaciona na reciprocidade entre estes campos e as relações sociais, não existindo \textbf{um} campo \textbf{que} se sobreponha aos outros no seio da sociedade.
A priorização de determinados campos somente ocorre quando são constituídos de fora para dentro.
No entanto, quando os campos adquirem autonomia são constituídos de dentro para fora, como enfatiza Bourdieu ( 2002, p. 9, tradução nossa ), em sua obra Campo Poder, Campo Intelectual : “ o campo intelectual, ao modo do campo magnético, constitui \textbf{um} sistema de linhas de força ”, \textbf{que} é constituído por “ [... ] agentes e instituições [ \textbf{que} ] estão em \textbf{uma} relação de forças \textbf{que} se opõem e se \textbf{agregam}, em sua estrutura específica, em \textbf{um} lugar e momento dados no tempo ”.
Os campos de saberes e práticas do ensino constituem -se, portanto, em unidades curriculares \textbf{que} substituirão as \textbf{disciplinas} tradicionais \textbf{atuais}.
Sua escolha se deu porque, por meio desses campos, é possível construir \textbf{um} terreno fértil para provocar a integração curricular entre os objetos de ensino \textbf{que} compõem \textbf{uma} área de conhecimento, a partir de \textbf{um} diálogo interdisciplinar em \textbf{um} trato contextualizado, mediante ao planejamento integrado dos professores de \textbf{uma} área e até mesmo entre duas ou mais áreas.
A escola é \textbf{um} \textbf{lócus} de formação \textbf{que} reúne \textbf{uma} diversidade de \textbf{sujeitos} \textbf{que} na construção coletiva do trabalho pedagógico deve visar práticas educativas \textbf{que} \textbf{fomentem} à formação humana integral.
Considerar as trajetórias da comunidade escolar e a relação entre os \textbf{sujeitos} é reconhecer as juventudes em sua inteireza, para potencializar, neste contexto de reforma, as reflexões necessárias, tendo em conta a capacidade crítica \textbf{que} há no ambiente escolar como espaço de emancipação, privilegiado e formativo das múltiplas dimensões humanas, \textbf{que} não estão centradas apenas no aluno, mas em todos os indivíduos \textbf{que} dão sentido e motivam a visão articulada dos campos de saberes das áreas de conhecimento, resultando avanços \textbf{que} ultrapassem o status quo, para a garantia da melhoria da aprendizagem e a permanência e conclusão do aluno com sucesso.
A Formação Geral Básica, por ter a intencionalidade de consolidar a aprendizagem, desenvolve, por meio das áreas de conhecimento, a aquisição de saberes, \textbf{que} no processo \textbf{desencadeado} pela ação pedagógico-docente requerem \textbf{que} o conhecimento esteja mutuamente articulado nas áreas de conhecimento e campos.
Dessa forma, as áreas dialogam com a concepção do documento curricular para \textbf{que} haja qualidade nos processos educativos, apresentando a sociedade as perspectivas \textbf{que} alteram a prática no ambiente escolar, expressa implicitamente pelas alterações da reforma e explicitamente na vigilância crítica, \textbf{que} aponta intervenções \textbf{que} se deseja alcançar.
Para melhor orientar a proposta para a escola, foi elaborado o Organizador Curricular da Formação Geral Básica, contendo a relação entre os princípios curriculares, as áreas de conhecimento, as \textbf{categorias} \textbf{conceituais} \textbf{que} fundamentam a área, as competências específicas de áreas, alinhados às habilidades.
A partir da associação desses elementos, são estabelecidas as possibilidades de objetos de conhecimento a serem consolidados.
A figura 11, a seguir, busca esquematizar essa relação : As áreas do conhecimento são articuladas pelos princípios, \textbf{que} são o aporte da etapa Ensino Médio, propondo e indicando caminhos \textbf{que} articulam as \textbf{categorias} das áreas, considerando o social, histórico, cultural e a tecnologia, alinhados aos processos e práticas de investigação \textbf{que} são conduzidos à construção coletiva de conhecimento dos \textbf{sujeitos} \textbf{que} fazem educação, escola-pedagógico-aluno-professor, \textbf{que} ocorre dentro e fora da escola no processo \textbf{dialógico}.
Ramos ( 2005, p. 116, \textbf{grifos} da autora ) afirma \textbf{que} : > [... ] o conhecimento não é de coisa, entidades, seres etc., mas sim das relações \textbf{que} se trata de descobrir, apreender no plano do pensamento.
São as apreensões assim elaboradas e formalizadas \textbf{que} constituem a \textbf{teoria} e os conceitos.
A Ciência é a parte do conhecimento mais bem \textbf{sistematizado} e deliberadamente expresso na forma de conceitos representativos das relações determinadas e apreendidas da realidade considerada.
O conhecimento é \textbf{uma} seção da realidade concreta ou a realidade concreta tematizada constitui os campos da ciência.
Dessa maneira, a Nucleação da Formação Geral Básica requer engajamento com o projeto histórico de ciência do conhecimento nos campos de saberes das áreas, para além da formação da consciência crítica, fazendo sentido ao \textbf{que} for apreendido pelos objetos de conhecimento.
O processo pedagógico, mediante a escuta dos jovens, contribui para a construção de \textbf{horizontes} a partir da mobilização dos interesses \textbf{que} se alargam no decorrer das aprendizagens, pois pensar o \textbf{sujeito} como elemento essencialmente ativo, \textbf{historicamente} constituído em seu espaço e tempo, \textbf{que} se percebe no território e estabelece \textbf{uma} práxis educativa para a construção de seus projetos de vida e suas lutas sociais, entre elas a educação, é \textbf{compromisso} do DCEPA – etapa ensino médio.
Assim, ao se discutir o processo de escuta das juventudes com propósito de abordar, de maneira qualificada, se retomará alguns pontos importantes \textbf{que} pressupõem a participação dos jovens e da comunidade escolar nas decisões dos processos educativos.
Essa participação tornou -se \textbf{uma} exigência social diante da realidade educacional do país, com a universalização do ensino, expõe -se a heterogeneidade das juventudes \textbf{que} estão presentes na escola, sobretudo, a pública.
No caso do ensino médio, a escola passa a perceber a diversidade das juventudes em suas \textbf{inúmeras} dimensões ; por isso, a formação humana integral procura atender essas expectativas da Resolução CNE/CEB nº 02/2012, enfatizando \textbf{que} os conhecimentos devem levar em conta sua organização na perspectiva da \textbf{totalidade} humana ( Art. 5º, inciso I ).
Vale ressaltar o percurso das referências \textbf{que} \textbf{embasaram} a participação ativa das juventudes nos processos de reformas educacionais, e \textbf{que}, a partir das conquistas democráticas, foram elaboradas.
A exemplo, em novembro de 1985, com a Lei nº 7.398, \textbf{que} dispõe sobre o direito de se organizar em entidades representativas e participação autônoma dos estudantes no ensino fundamental e médio, o Grêmio Estudantil, \textbf{configurou} -se como \textbf{um} grande avanço para participação dos jovens nos estabelecimentos de ensino.
Nesse sentido, a Constituição Federal de 1988 estabelece a gestão democrática como caminho importante a ser considerado, pois a \textbf{democratização} da gestão é \textbf{uma} premissa para a qualidade da educação, \textbf{uma} vez \textbf{que} criou possibilidades de a escola potencializar a participação efetiva dos diferentes \textbf{atores} da comunidade escolar ( pais, professores, estudantes e funcionários ) em todos os âmbitos da organização educativa.
No ano de 1990, \textbf{um} marco histórico foi o ECA ( Lei nº 8.069/1990 ), sendo considerado os seus direitos civis, humanos e sociais, dando ênfase à liberdade de expressão, à opinião e à participação política na educação.
Consolida o direito do aluno a ser respeitado pelos educadores e o direito dos estudantes à organização e participação em entidades estudantis.
A LDB em 1996 ( Lei 9.394/1996 ) reforçou a gestão democrática mediante aos princípios da participação das comunidades locais e escolar nos conselhos escolares.
O Estatuto das Juventudes em 2013 ( Lei nº 12.852/2013 ) vem como referência para consolidar os direitos de jovens de 15 a 29 anos.
Em seu artigo 2º, abordam -se os princípios \textbf{que} contribuem para \textbf{fomentar} a escuta dos estudantes na escola.
Tem -se a promoção da autonomia, da participação política, social e do diálogo, assim como o respeito à identidade e à diversidade individual e coletiva das juventudes.
Ainda de acordo com o Estatuto da Juventude, em suas Diretrizes Gerais, no artigo 3{\EmojiFont °}, os agentes públicos devem incentivar a ampla participação juvenil em formulação, implementação e avaliação de políticas públicas.
No Plano Nacional de Educação ( PNE ) – Lei nº 13.005/2014 foram garantidos o fortalecimento dos grêmios estudantis e sua articulação com os conselhos escolares, bem como a sua constituição, sendo instrumentos de participação e fiscalização da gestão escolar.
Dessa maneira foram garantidos a participação dos estudantes na construção dos Projetos Políticos-Pedagógicos ( PPP ), currículos escolares, planos da gestão escolar e regimentos escolares.
Na perspectiva da concepção humana integral, \textbf{que} compreende o desenvolvimento das juventudes em todas as suas dimensões e sua \textbf{centralidade} em todos os aspectos da proposta pedagógica da escola, a escuta é \textbf{um} processo \textbf{que} se realiza não somente com o \textbf{educando}, mas deve ser de todos os segmentos \textbf{que} fazem parte da escola, desde a concepção do PPP até a definição da oferta no âmbito escolar, na qual a coletividade possa defender e propor organizações curriculares e projetos pedagógicos \textbf{que} expressem as escolhas, os interesses e as realidades \textbf{que} são concebidas por meio da participação dos \textbf{sujeitos} \textbf{que} compõem a escola.
Por essas razões, a escuta desses estudantes é \textbf{uma} prática \textbf{imprescindível} na organização do trabalho pedagógico da escola, a qual torna -se \textbf{um} ato educativo no qual o aluno externaliza os seus sentimentos, seus anseios e suas \textbf{vontades} como \textbf{sujeito} do processo de \textbf{ensino-aprendizagem}.
Ouvir os estudantes significa criar oportunidade para \textbf{que} possam compartilhar opiniões sobre diferentes assuntos, desde os mais corriqueiros, como a infraestrutura da escola e as atividades em sala de aula, até os mais complexos, como mudanças no currículo e na organização escolar.
Esse processo de \textbf{atentar} às necessidades dos jovens, às suas vozes no âmbito escolar, precisa ser \textbf{um} movimento contínuo, \textbf{que} por meio das falas, são firmadas expectativas, anseios e \textbf{compromisso} de quem faz parte da comunidade escolar, pois o jovem, quando se \textbf{vê} como partícipe do processo educativo, cria expectativas sobre suas decisões e tornando -se \textbf{corresponsável} por elas ; assim acontece com os demais \textbf{sujeitos} do processo escolar.
Por isso é importante \textbf{institucionalizar} esses espaços escolares, como grêmio estudantil, conselhos de classe, conselho escolar, sendo esses importantes instrumentos de diálogo e de participação nas decisões da escola.
Assim, a escola deve organizar discussões conjuntas com a comunidade escolar, podendo utilizar -se de reuniões presenciais, assim como questionários, pesquisas, consultas online, com intencionalidade de mudanças no currículo e na sua organização.
Após esse primeiro momento, a escola precisa criar \textbf{um} processo de \textbf{sistematização} em \textbf{que} deverá constar a porcentagem dos interesses dos estudantes.
Esses dados devem estar obrigatoriamente no PPP das escolas, e disponíveis por meio de registro dos anseios e sugestões estudantis, tendo em vista a promoção do protagonismo juvenil, partindo de \textbf{um} diagnóstico de interesses e expectativas de aprendizagem \textbf{que} deverão ser desenvolvidas por meio da proposta curricular e, ao mesmo tempo, apoiar suas escolhas em relação às itinerâncias, \textbf{fomentando} a participação cidadã desse jovem.
Nesse sentido, a escola precisa \textbf{institucionalizar} \textbf{um} processo de escuta amplo \textbf{que} considere seu projeto político-pedagógico e culmine no diagnóstico dos interesses e expectativas de aprendizagem dos estudantes no ensino médio.
É importante pautar \textbf{que} o diagnóstico dos interesses e expectativas de aprendizagem são elementos \textbf{que} contribuem para o aluno pensar o seu projeto de vida, \textbf{uma} vez \textbf{que} este favorece o desenvolvimento dos jovens no \textbf{que} \textbf{diz} respeito às suas identidades e sua relação de pertencimento com a comunidade escolar.
A escuta adquire \textbf{um} maior significado quando associada ao projeto de vida do aluno, tendo em vista \textbf{que} este trata da construção das escolhas de vida, dos desejos de mudança, da realização de propósitos e da projeção de futuro.
Para melhor orientar esse processo, apresentam -se alguns procedimentos necessários para as instituições escolares efetivarem as suas escutas \textbf{que} estão listados a seguir, no quadro 04, com base nas contribuições de Padilha ( 2008 ) : [ 49 ] : Sobre esta segunda nucleação, o item 4.2 A nucleação da formação para o mundo do trabalho, abordará essa discussão, a seguir.

% matched lemmas: agregar, amparar, aproximar, atentar, ator, atual, autor, basilar, categoria, centralidade, compromisso, conceitual, configurar, conseguinte, contemporâneo, contextualização, contudo, corresponsável, democratização, desencadear, despertar, dialeticamente, dialético, dialógico, disciplina, dizer, educar, embasar, enquanto, ensino-aprendizagem, equidade, fomentar, fragmentação, grifo, hegemônico, historicamente, horizonte, imprescindível, influenciar, influência, institucionalizar, inter-relação, interdisciplinaridade, inúmero, lado, literário, lócus, norteador, perder, permear, preciso, qualificar, que, remeter, sistematizar, sistematização, sujeito, supor, século, sócio-histórico, teoria, teórico-metodológico, totalidade, um, ver, viver, vontade
\end{textsample}
